This report presented a \MATLAB{} implementation of argument-of-latitude
optimisation for a six-satellite LEO constellation. The key contribution
is a reduction of the $N$-dimensional scheduling problem to a scalar
optimisation over a single reference phase~$\uref$, exploiting the
equal-spacing constraint. The two-burn phasing cost is computed
analytically from orbital mechanics, and \texttt{fminbnd} finds the
optimal reference phase on $[0,2\pi)$.

Two strategies were implemented and compared on a 550~km-altitude test
case. The \textbf{weighted strategy} ($W_1=1$, $W_2=10$) produces a
more balanced remaining-fuel distribution at a slightly higher total $\dv$.
The \textbf{lexicographic strategy} strictly prioritises fuel economy in
Stage~1 and then improves balance in Stage~2 within a tight tolerance,
yielding results comparable to the weighted approach for the test case.

\subsection{Limitations and Possible Extensions}

\begin{description}
  \item[Multi-revolution phasing.]
    Each transfer is currently fixed to one phasing revolution. Allowing
    $k$ revolutions reduces $\dv$ by roughly $1/k$ at the cost of a
    longer transfer time---a worthwhile trade-off when large angular
    corrections are required.

  \item[Feasibility check.]
    When the computed $\dv_i$ exceeds the satellite's budget $f_i$,
    the solver returns negative remaining fuel without a warning.
    A feasibility guard should be added.

  \item[Global optimality.]
    \texttt{fminbnd} is a local search; for non-unimodal objectives,
    a multi-start sweep over $\uref$ would strengthen the global-optimality
    guarantee.

  \item[Perturbation model.]
    Neglecting $J_2$, atmospheric drag, and solar pressure during the
    transfer leads to optimistic $\dv$ estimates.

  \item[Non-circular orbits.]
    Generalising to elliptic reference orbits would require a Hohmann
    transfer followed by a phasing arc, substantially increasing the
    complexity of the $\dv$ calculation.
\end{description}
